\plainsection{表 2}

我们把每个基本权与相应单表示的维数联系起来，该维数由§9的定理 2 计算得到。

\[
\begin{array}{c c c} \mathrm{A} _ {l} (l \geq 1) & & \mathrm{B} _ {l} (l \geq 2) \\ \varpi_ {r} (1 \leq r \leq l) \quad \binom {l + 1} {r} & & \varpi_ {r} (1 \leq r \leq l - 1) \quad \binom {2 l + 1} {r} \\ & & \varpi_ {l} \quad 2 ^ {l} \end{array}
\]

\[
\begin{array}{c c c} \mathrm{C} _ {l} (l \geq 2) & \mathrm{D} _ {l} (l \geq 2) \\ \varpi_ {r} (1 \leq r \leq l) \quad \binom {2 l} {r} - \binom {2 l} {r - 2} & \varpi_ {r} (1 \leq r \leq l - 2) & \binom {2 l} {r} \\ & \varpi_ {l - 1} & 2 ^ {l - 1} \\ & \varpi_ {l} & 2 ^ {l - 1} \end{array}
\]

\[
\mathrm{E} _ {6}
\]

\[
\varpi_ {1} 27 = 3 ^ {3}
\]

\[
\mathrm{E} _ {7}
\]

\[
\varpi_ {1} 133 = 7.19
\]

\[
\varpi_ {2} 78 = 2.3.13
\]

\[
\varpi_ {2} 912 = 2 ^ {4}.3.19
\]

\[
\varpi_ {3} 351 = 3 ^ {3}.13
\]

\[
\varpi_ {3} 8645 = 5.7.13.19
\]

\[
\varpi_ {4} 2925 = 3 ^ {2}.5 ^ {2}.13
\]

\[
\varpi_ {4} 365750 = 2.5 ^ {3}.7.11 ^ {2}.17.23.31
\]

\[
\varpi_ {5} 351 = 3 ^ {3}.13
\]

\[
\varpi_ {5} 27664 = 2 ^ {4}.7.13.19
\]

\[
\varpi_ {6} 27 = 3 ^ {3}
\]

\[
\varpi_ {6} 1539 = 3 ^ {4}.19
\]

\[
\varpi_ {7} 56 = 2 ^ {3}.7
\]

\[
\mathrm{E} _ {8}
\]

\[
\varpi_ {8} 248 = 2 ^ {3}.31
\]

\[
\mathrm{F} _ {4}
\]

\[
\varpi_ {1} 3875 = 5 ^ {3}.31
\]

\[
\varpi_ {1} 52 = 2 ^ {2}.13
\]

\[
\varpi_ {2} 147250 = 2.5 ^ {3}.19.31
\]

\[
\varpi_ {2} 1274 = 2.7 ^ {2}.13
\]

\[
\varpi_ {3} 6696000 = 2 ^ {6}.3 ^ {3}.5 ^ {3}.31
\]

\[
\varpi_ {3} 273 = 3.7.13
\]

\[
\varpi_ {4} 6899079264 = 2 ^ {5}.3.7 ^ {2}.11 ^ {2}.17.23.31
\]

\[
\varpi_ {4} 26 = 2.13
\]

\[
\varpi_ {5} 146325270 = 2.3.5.7 ^ {2}.13 ^ {2}.19.31
\]

\[
\varpi_ {6} 2450240 = 2 ^ {6}.5.13.19.31
\]

\[
\varpi_ {7} 30380 = 2 ^ {2}.5.7 ^ {2}.31
\]

\[
\mathrm{G} _ {2}
\]

\[
\varpi_ {1} 7
\]

\[
\varpi_ {2} 14 = 2.7
\]

\specialchapter{习题}

设基域 \(k\) 的特征为零。

除非另有明确说明，李代数均假定为有限维的。

\exercisesection{\S~1}

记 \(\mathfrak{s}\) 为李代数 \(\mathfrak{sl}(2,k)\)。

\begin{enumerate}
\def\labelenumi{\arabic{enumi})}
\tightlist
\item
  设 \(\mathrm{U}\) 为 \(\mathfrak{s}\) 的包络代数。证明元素
\end{enumerate}

\[
\mathrm{C} = H ^ {2} - 2 (X _ {+} X _ {-} + X _ {-} X _ {+}) = H ^ {2} + 2 H - 4 X _ {-} X _ {+}
\]

属于 U 的中心，并且它在与 \(V(m)\) 相关的表示中的像是比值为 \(m(m+2)\) 的伸缩变换。

\begin{enumerate}
\def\labelenumi{\arabic{enumi})}
\setcounter{enumi}{1}
\tightlist
\item
  设 \(\lambda \in k\)，设 \(Z(\lambda)\) 为具有基 \((e_n)\) 的向量空间，其中 \(n = 0,1,\ldots\)，并设 \(X_{+}, X_{-}, H\) 为由命题 1 的公式 (2) 定义的 \(Z(\lambda)\) 的自同态。
\end{enumerate}

\begin{enumerate}
\def\labelenumi{\alph{enumi})}
\item
  验证这给出一个\(\mathfrak{s}\)-模结构，作用于\(Z(\lambda)\)。
\item
  假设 \(\lambda\) 不是一个 \(\geq 0\) 的整数。证明 \(\mathfrak{s}\)-模 \(Z(\lambda)\) 是单的。
\item
  假设 \(\lambda\) 是一个 \(\geq 0\) 的整数。令 \(Z'\) 为 \(Z(\lambda)\) 的子空间，由 \(e_{n}\)（\(n > \lambda\)）生成。证明 \(Z'\) 是 \(Z(\lambda)\) 的 s-子模，同构于 \(Z(-\lambda - 2)\)，且商 \(Z(\lambda)/Z'\) 同构于单 s-模 \(V(\lambda)\)。证明 \(Z(\lambda)\) 的 s-子模只有 0、\(Z'\) 和 \(Z(\lambda)\)。
\end{enumerate}

\begin{enumerate}
\def\labelenumi{\arabic{enumi})}
\setcounter{enumi}{2}
\tightlist
\item
  设 \(\mathrm{E}\) 为具有基 \((e_n)_{n\in \mathbf{Z}}\) 的向量空间。令
\end{enumerate}

\[
a (n) = a _ {0} + a _ {1} n, b (n) = b _ {0} + b _ {1} n, c (n) = c _ {0} + c _ {1} n
\]

为系数属于 k 的三个仿射函数。按下列公式定义 E 的自同态 \(X_{+}, X_{-}\)、H：

\[
X _ {+} e _ {n} = a (n) e _ {n - 1}, X _ {-} e _ {n} = b (n) e _ {n + 1}, H e _ {n} = c (n) e _ {n}.
\]

当且仅当满足下列条件时，这给出 E 上的一个 \(\mathfrak{s}\)-模结构：

\[
a _ {1} b _ {1} = 1, c _ {1} = - 2, c _ {0} = - a _ {0} b _ {1} - a _ {1} b _ {0}.
\]

此 \(\mathfrak{s}\)-模在什么条件下是单的？

¶ 4) 设 g 为李代数。若 g-模 E 是有限维 g-子模的并，则称其为局部有限；等价地，E 的每个有限型（作为 U(g)-模的）g-子模都是有限维的。

\begin{enumerate}
\def\labelenumi{\alph{enumi})}
\item
  设 \(0 \rightarrow E' \rightarrow E \rightarrow E'' \rightarrow 0\) 是 g-模的正合序列。证明，如果 \(E'\) 和 \(E''\) 局部有限，则 E 也是局部有限（归约到 E 是有限型的情形，并利用 \(\mathrm{U}(\mathfrak{g})\) 是诺特环这一事实，参见~第 I 章，§2，节 6）。
\item
  假设 g 是半单的。证明 E 局部有限当且仅当 E 是有限维单 g-模的直和。
\item
  假设 \(\mathfrak{g} = \mathfrak{s}\)。证明 E 局部有限当且仅当满足以下条件：
\end{enumerate}

\(c_{1})\) E 有一个由 H 的特征向量组成的基，

\(c_{2}\) ) 自同态 \(X_{+\mathrm{E}}\) 和 \(X_{-\mathrm{E}}\) 局部幂零。

（先证明：如果 E 满足 \(c_{1}\) 和 \(c_{2}\)，且 E 不等于 0，则它包含一个具有整数权 \(m \geq 0\) 的本原元 e；再证明 \(X^{m+1}e = 0\)，从而 E 包含 \(V(m)\)。最后将 a) 应用于 E 的最大局部有限子模 \(E'\)。）

\begin{enumerate}
\def\labelenumi{\arabic{enumi})}
\setcounter{enumi}{4}
\tightlist
\item
  设 E 是一个局部有限的 \(\mathfrak{s}\)-模（习题 4）。对所有 \(m \geq 0\)，记 \(\mathrm{L}(m)\) 为 \(\mathfrak{s}\)-同态从 \(\mathrm{V}(m)\) 到 E 所成的向量空间。
\end{enumerate}

\begin{enumerate}
\def\labelenumi{\alph{enumi})}
\item
  定义从 E 到 \(\bigoplus_{m} \mathrm{L}_{m} \otimes \mathrm{V}(m)\) 的一个同构。
\item
  记 \(\Phi_{m}\) 为第 3 节备注 3 中定义的 \(\mathrm{V}(m)\) 上的不变双线性型。对所有 \(m \in N\)，令 \(b_{m}\) 为 \(L_{m}\) 上的双线性型；令 b 为 E 上的双线性型，它在 a) 中的同构下对应于双线性型 \(b_{m} \otimes \Phi_{m}\) 的直和。证明 b 是不变的，并且 E 上的每个不变双线性型都以唯一方式由此得到。b 是对称的（分别为交错的）当且仅当 m 为偶数时的 \(b_{m}\) 是对称的（分别为交错的），而 m 为奇数时的 \(b_{m}\) 是交错的（分别为对称的）。b 非退化当且仅当 \(b_{m}\) 非退化。
\item
  若 E 是有限维的，证明 E 是单生成的（作为 \(\mathrm{U}(\mathfrak{s})\)-模）当且仅当 \(\dim L_{m} \leq m + 1\) 对所有 \(m \geq 0\) 都成立。
\end{enumerate}

\begin{enumerate}
\def\labelenumi{\arabic{enumi})}
\setcounter{enumi}{5}
\tightlist
\item
  若 E 是有限维的 \(\mathfrak{s}\)-模，且 \(n\) 为一个整数，记 \(a_{n}\) 为 \(H_{\mathrm{E}}\) 的特征值为 \(n\) 的特征子空间的维数。记 \(c_{\mathrm{E}}(\mathrm{T})\) 为 \(\mathbf{Z}[\mathrm{T},\mathrm{T}^{-1}]\) 中由 \(c_{\mathrm{E}}(\mathrm{T}) = \sum_{n\in \mathbf{Z}}a_n\mathrm{T}^n\) 定义的元素。
\end{enumerate}

\begin{enumerate}
\def\labelenumi{\alph{enumi})}
\tightlist
\item
  按习题 5 定义 \(\mathrm{L}_m\)。证明
\end{enumerate}

\(\dim \mathrm{L}_m = a_m - a_{m + 2}\) 对每个整数 \(m\geq 0\)

推出，\(c_{\mathrm{E}}(\mathrm{T}) = c_{\mathrm{E}'}(\mathrm{T})\) 当且仅当 E 与 \(E'\) 同构。利用第 VII 章 §3 的习题 18 e) 重新得到这一结果。

\begin{enumerate}
\def\labelenumi{\alph{enumi})}
\setcounter{enumi}{1}
\item
  证明 \(c_{E\oplus F} = c_{E} + c_{F}\) 和 \(c_{E\otimes F} = c_{E}.c_{F}\)。
\item
  有 \(c_{\mathrm{V}(m)}(\mathrm{T}) = (\mathrm{T}^{m + 1} - \mathrm{T}^{-m - 1}) / (\mathrm{T} - \mathrm{T}^{-1})\)。
\end{enumerate}

\(d)\) 由 \(a\)、\(b\)、\(c)\) 推出：如果 \(m \geq m' \geq 0\)，则 \(\mathfrak{s}\)-模 \(\mathrm{V}(m) \otimes \mathrm{V}(m')\) 同构于

\[
\mathrm{V} (m + m ^ {\prime}) \oplus \mathrm{V} (m + m ^ {\prime} - 2) \oplus \mathrm{V} (m + m ^ {\prime} - 4) \oplus \dots \oplus \mathrm{V} (m - m ^ {\prime}).
\]

若 \(m\) 为 \(\geq 0\)，则 \(\mathfrak{s}\)-模 \(\mathbf{S}^2\mathrm{V}(m)\) 同构于

\[
\mathrm{V} (2 m) \oplus \mathrm{V} (2 m - 4) \oplus \mathrm{V} (2 m - 8) \oplus \dots \oplus \left\{ \begin{array}{l l} \mathrm{V} (0) & \text {if m is even} \\ \mathrm{V} (2) & \text {if m is odd} \end{array} \right.
\]

且 \(\mathfrak{s}\)-模 \(\bigwedge^2\mathrm{V}(m)\) 同构于

\[
\mathrm{V} (2 m - 2) \oplus \mathrm{V} (2 m - 6) \oplus \mathrm{V} (2 m - 10) \oplus \dots \oplus \left\{ \begin{array}{l l} \mathrm{V} (2) & \text {if m is even} \\ \mathrm{V} (0) & \text {if m is odd.} \end{array} \right.
\]

\begin{enumerate}
\def\labelenumi{\arabic{enumi})}
\setcounter{enumi}{6}
\item
  设 E 是有限维 s-模。证明 E 的对偶同构于 E。
\item
  证明 \(\mathfrak{s}\)-模 \(V(m)\) 可实现为两个变量中的齐次多项式 \(f(u,v)\)（次数为 \(m\)）所成的空间，其中算子 \(X_{+}, X_{-}\) 和 \(H\) 由下式给出：
\end{enumerate}

\[
X _ {+} f = u \frac {\partial f}{\partial v}, X _ {-} f = - v \frac {\partial f}{\partial u}, H f = u \frac {\partial f}{\partial u} - v \frac {\partial f}{\partial v}.
\]

¶9) 考虑 \(\mathfrak{s}\) 通过伴随表示在其包络代数 U 和其对称代数 \(\mathbf{S} = \bigoplus \mathbf{S}^n\) 上的作用。

\begin{enumerate}
\def\labelenumi{\alph{enumi})}
\tightlist
\item
  确定 \(\mathfrak{s}\)-模 \(\mathbf{S}^n\) 的权。推出以下 \(\mathfrak{s}\)-模同构：
\end{enumerate}

\[
\begin{array}{l l} \mathbf {S} ^ {n} \longrightarrow \mathrm{V} (2 n) \oplus \mathrm{V} (2 n - 4) \oplus \mathrm{V} (2 n - 8) \oplus \dots \oplus \mathrm{V} (0) & n \text {even} \\ \mathbf {S} ^ {n} \longrightarrow \mathrm{V} (2 n) \oplus \mathrm{V} (2 n - 4) \oplus \mathrm{V} (2 n - 8) \oplus \dots \oplus \mathrm{V} (2) & n \text {odd}. \end{array}
\]

特别地，\(\mathbf{S}^n\) 中在 \(\mathfrak{s}\) 下不变的元素构成一个空间；当 \(n\) 为偶数（分别为奇数）时，其维数为 1（分别为 0）。

\begin{enumerate}
\def\labelenumi{\alph{enumi})}
\setcounter{enumi}{1}
\item
  证明由在 s 下不变的元素组成的 S 的子代数是多项式代数 \(k[\Gamma]\)，其中 \(\Gamma = H^{2} - 4X_{-}X_{+}\)。模 \(S^{n}/\Gamma.S^{n-2}\) 同构于 \(\mathrm{V}(2n)\)。
\item
  证明 U 的中心是多项式代数 \(k[\mathbf{C}]\)，其中 C 是习题 1 中定义的元素（利用 b) 和 Poincaré-Birkhoff-Witt 定理）。
\end{enumerate}

\begin{enumerate}
\def\labelenumi{\arabic{enumi})}
\setcounter{enumi}{9}
\tightlist
\item
  设 \(m\) 为一个 \(> 0\) 的整数，\(\mathbf{S}\) 为分次代数 \(k[\mathrm{X}_1, \ldots, \mathrm{X}_m]\)，且 \(a_1, \ldots, a_m\) 为 \(k^*\) 中的元素。设 \(\Phi = \sum_{i,j=1}^{m} a_{ij} X_i X_j\) 为一个二次型；置 \(\mathrm{D}_i = \frac{\partial}{\partial X_i}\)，并置 \(\mathrm{D} = \sum_{i,j=1}^{m} b_{ij} \mathrm{D}_i \mathrm{D}_j\)，其中 \((b_{ij})\) 是 \((a_{ij})\) 的逆矩阵。
\end{enumerate}

\begin{enumerate}
\def\labelenumi{\alph{enumi})}
\tightlist
\item
  证明存在唯一的 \(\mathfrak{s}\)-模结构，作用在 \(\mathbf{S}\) 上，使得 \(X_{+}f = \frac{1}{2}\mathrm{D}(f), X_{-}f = \frac{1}{2}\Phi f\)，并且 \(Hf = -(m / 2 + n)f\)（当 \(f\) 是次数为 \(n\) 的齐次元素时）。b) 置 \(\mathrm{A}^n = \mathbf{S}^n \cap \mathrm{Ker}(\mathrm{D})\)。证明 \(\mathbf{S}^n\) 是 \(\Phi^p\)、\(\mathrm{A}^q\)（其中 \(2p + q = n\)）这些子空间的直和。推出恒等式
\end{enumerate}

\[
\sum_ {n = 0} ^ {\infty} \dim (\mathrm{A} ^ {n}) \mathrm{T} ^ {n} = (1 + \mathrm{T}) / (1 - \mathrm{T}) ^ {m - 1}.
\]

\begin{enumerate}
\def\labelenumi{\alph{enumi})}
\setcounter{enumi}{2}
\item
  若 \(f\) 是 \(\mathbf{A}^n\) 的非零元素，则 \(\varPhi^p f\)（\(p \geq 0\)）构成一个单的 \(\mathfrak{s}\)-子模的基；该子模属于 \(\mathbf{S}\)，并同构于习题 2 中的模 \(Z\left(-\frac{m}{2} - n\right)\)。
\item
  具体写出 \(m = 1\) 和 \(m = 2\) 的情形。利用 \(m = 3\) 的情形重新得到习题 9 的结果。
\end{enumerate}

\begin{enumerate}
\def\labelenumi{\arabic{enumi})}
\setcounter{enumi}{10}
\tightlist
\item
  假设 \(k\) 是代数闭的。设 \(\mathfrak{g}\) 为李代数，\(V\) 为单 \(\mathfrak{g}\)-模，且 \(D\) 为 \(\mathfrak{g}\)-模 \(V\) 的自同态域。证明，如果 \(k\) 是不可数的，\(^7\) 则 \(D = k\)。（否则，\(D\) 将包含一个同构于 \(k(X)\) 的子域，其中 \(X\) 是不定元，于是有 \(\dim_k D > \aleph_0\)，从而也有 \(\dim_k V > \aleph_0\)，这与 \(V\) 是单生成的 \(U(\mathfrak{g})\)-模相矛盾。）
\end{enumerate}

\footnotetext[7]{结论即使 k 可数时仍然成立，参见 D. QUILLEN，《关于包络代数上一个单模的自同态环》，Proc. Amer. Math. Soc., Vol. XXI (1969), pp. 171-172。}

¶ 12) 设 \(q \in k\)，且 \(W\) 为具有基 \((e_0, e_1, e_2, \ldots)\) 的向量空间。

\begin{enumerate}
\def\labelenumi{\alph{enumi})}
\tightlist
\item
  证明存在唯一的表示 \(\rho_{q}\)，它是 \(\mathfrak{s}\) 上的表示并作用在 W 上，使得
\end{enumerate}

\[
\begin{array}{r l} & {\rho_ {q} (H) e _ {n} = 2 e _ {n + 1}, \quad \rho_ {q} (X _ {+}) e _ {n} = (\frac {1}{2} \rho_ {q} (H) - 1) ^ {n} e _ {0},} \\ & {\rho_ {q} (X _ {-}) e _ {n} = (\frac {1}{2} \rho_ {q} (H) + 1) ^ {n} (- q e _ {0} + e _ {1} + e _ {2}).} \end{array}
\]

有 \(\rho_{q}(\mathrm{C}) = 4q\)，其中 \(C = H^{2} + 2H - 4X_{-}X_{+}\)（参见~习题 1）。表示 \(\rho_{q}\) 是单的。元素 \(x \in s\) 使 \(\rho_{q}(x)\) 具有特征值的，正是 \(X_{+}\) 的倍数。自同态 \(\rho_{q}(X_{+})\) 以重数 1 具有特征值 1。

\begin{enumerate}
\def\labelenumi{\alph{enumi})}
\setcounter{enumi}{1}
\tightlist
\item
  设 \(\rho\) 是 \(\mathfrak{s}\) 的一个单表示，满足 \(\rho(\mathrm{C}) = 4q\)，且 \(\rho(X_{+})\) 具有特征值 1。证明 \(\rho\) 等价于 \(\rho_q\)。
\end{enumerate}

¶ 13) 假设 \(k = \mathbf{C}\)。置 \(C = H^2 + 2H - 4X_-X_+\)，参见~习题 1。一个表示 \(\rho\)（其对象为 \(\mathfrak{s} = \mathfrak{sl}(2, \mathbf{C})\)）被称为 \(H\)-可对角化的，如果 \(\rho\) 的底层空间有一个由 \(\rho(H)\) 的特征向量组成的基。令 \(q \in \mathbf{C}\)，\(v \in \mathbf{Q}/\mathbf{Z}\)。

\begin{enumerate}
\def\labelenumi{\alph{enumi})}
\tightlist
\item
  设 S 为具有基 \((e_{w})_{w\in\mathbf{C}}\) 的向量空间，该基由 C 的元素索引。令 \(S_{v}=\sum_{w\in v}Ce_{w}\)。存在唯一的表示 \(\rho_{v,q}\)，它是 s 在 \(S_{v}\) 上的表示，使得
\end{enumerate}

\[
\rho_ {v, q} (X _ {+}) e _ {w} = (q - w ^ {2} - w) ^ {1 / 2} e _ {w + 1}
\]

\[
\begin{array}{c} \rho_ {v, q} (X _ {-}) e _ {w} = (q - w ^ {2} + w) ^ {1 / 2} e _ {w - 1} \\ \rho_ {v, q} (H) e _ {w} = 2 w e _ {w}. \end{array}
\]

（约定，对于所有 \(z \in \mathbf{C}\)，\(z^{1/2}\) 是 \(z\) 的平方根，其幅角属于 \(\{0, \pi\}\)。）记 \(\mathrm{S}_{v,q}\) 为一个 \(\mathfrak{s}\)-模，它由 \(\rho_{v,q}\) 定义。有 \(\rho_{v,q}(\mathrm{C}) = 4q\)。

\begin{enumerate}
\def\labelenumi{\alph{enumi})}
\setcounter{enumi}{1}
\item
  若 \(q \neq u^2 + u\) 对所有 \(u \in v\) 都成立，则 \(S_{v,q}\) 是单的。
\item
  假设 \(2v \neq 0\)，且 q 形如 \(u^{2} + u\)，其中 \(u \in v, u \geq 0\)（这唯一确定 u）。令 \(S_{v,q}^{-}\)（分别为 \(S_{v,q}^{+}\)）为 \(S_{v}\) 中由 \(e_{w}\)（满足 \(w \leq u\)；分别为 w \textgreater{} u）生成的向量子空间。则 \(S_{v,q}^{-}\) 和 \(S_{v,q}^{+}\) 是 \(S_{v,q}\) 的单 s-子模。
\item
  假设 2v = 0，且 q 形如 \(u^{2} + u\)，其中 \(u \in v, u \geq 0\)（这唯一确定 u）。令 \(S_{v,q}^{-}\)（分别为 \(S_{v,q}^{0}, S_{v,q}^{+}\)）为 \(S_{v}\) 中由 \(e_{w}\)（满足 w \textless{} -u；分别为 \(-u \leq w \leq u, w > u\)）生成的向量子空间。则 \(S_{v,q}^{-}, S_{v,q}^{0}\) 和 \(S_{v,q}^{+}\) 是 \(S_{v,q}\) 的单 s-子模。
\item
  假设 \(v = -\frac{1}{2} + Z\) 且 \(q = -\frac{1}{4}\)。令 \(S_{-1/2,-1/4}^{-}\)（分别为 \(S_{-1/2,-1/4}^{+}\)）为 \(S_{-1/2}\) 中由 \(e_{w}\)（满足 \(w \leq -\frac{1}{2}\)；分别为 \(w > -\frac{1}{2}\)）生成的向量子空间。则 \(S_{-1/2,-1/4}^{-}\) 和 \(S_{-1/2,-1/4}^{+}\) 是 \(S_{-1/2,-1/4}\) 的单 s-子模。
\end{enumerate}

\(f)\) 记 \(\rho_{v,q}^{\pm},\rho_{v,q}^{0}\) 为分别与 \(\mathrm{S}_{v,q}^{\pm},\mathrm{S}_{v,q}^{0}\) 对应的表示。在情形 \(b)\) 中，元素 \(x\in \mathfrak{s}\) 若使 \(\rho_{v,q}^{-}(x)\) 具有特征值，则这些元素正是 \(\mathbf{CH}\) 中的元素。在情形 \(c)\)、\(d)\)、\(e)\) 中，元素 \(x\in \mathfrak{s}\) 若使 \(\rho_{v,q}^{-}(x)\) 具有特征值，则这些元素正是 \(\mathbf{CH} + \mathbf{CX}_+\) 中的元素；此外，若 \(x\) 是幂零的（因而与 \(X_{+}\) 成比例）且非零，则 \(\rho_{v,q}^{-}(x)\) 的唯一特征值为 0，且其重数为 1；另一方面，若 \(x\) 是半单的，则 \(\rho_{v,q}^{-}\) 的底层空间有一个由 \(\rho_{v,q}^{-}(x)\) 的特征向量组成的基。对于 \(\rho_{v,q}^{+}\) 有类似结果，只需将 \(X_{+}\) 换成 \(X_{-}\)。

\(g)\) 设 \(V\) 是单的 \(\mathfrak{s}\)-模，\(\rho\) 为相应表示。则 \(\rho(C)\) 是伸缩变换（利用习题 11）。假设 \(\rho(C) = 4q\)。证明，如果 \(\rho\) 是 \(H\)-可对角化的，则 \(V\) 同构于模 \(S_{v,q}, S_{v,q}^{\pm}, S_{v,q}^{0}\)，该模在 \(b), c), d), e)\) 中被考虑。此外，\(\rho\) 是 \(H\)-可对角化的当且仅当 \(\rho(H)\) 具有特征值；只要 \(\rho(X_{+})\) 具有特征值 \(0\) 即可。

¶14) 记号沿用上一个习题。记 \(B_{q}\) 为 \(U(\mathfrak{s})\) 关于由 C - 4q 生成的双边理想的商，并记 \(u \mapsto u^{\bullet}\) 为 \(U(\mathfrak{s}) \to B_{q}\) 的典则映射。将习题 12 和 13 的表示看作 \(B_{q}\) 的表示。

\begin{enumerate}
\def\labelenumi{\alph{enumi})}
\tightlist
\item
  \(B_{q}\) 的每个元素都能唯一表示为
\end{enumerate}

\[
\sum_ {r \geq 0} X _ {+} ^ {\bullet r} p _ {r} (H ^ {\bullet}) + \sum_ {s > 0} q _ {s} (H ^ {\bullet}) X _ {-} ^ {\bullet s},
\]

其中 \(p_{r}\) 和 \(q_{s}\) 是多项式。若 \(B_{q}\) 的两个元素生成同一个左理想，则它们成比例。

\begin{enumerate}
\def\labelenumi{\alph{enumi})}
\setcounter{enumi}{1}
\item
  考虑习题 13 的情形 b)。设 a 是 \(S_{v,q}\) 的非零元素。如果 a 是 \(\rho_{v,q}(H)\) 的特征向量，特征值为 \(\lambda\)，则 a 在 \(B_{q}\) 中的湮没子是由 \(H^{\bullet}-\lambda\) 生成的左理想；如果 a 不是 \(\rho_{v,q}(H)\) 的特征向量，则其湮没子是非单生成左理想。
\item
  考虑习题 13 的表示 \(\rho_{v,q}^{\pm}\) 的情形。如果 a 是这种表示的空间中的非零元素，则其在 \(B_{q}\) 中的湮没子是非单生成左理想。
\item
  考虑习题 12 的表示 \(\rho_{q}\)。\(e_{0}\) 在 \(B_{q}\) 中的湮没子由 \(X_{+}^{\bullet}-1\) 生成。如果 \(a\in W\) 不与 \(e_{0}\) 成比例，则其湮没子是非单生成左理想。
\item
  每个 \(\mathfrak{s}\) 的自同构都延拓为 \(\mathrm{U}(\mathfrak{s})\) 的自同构，并通过取商定义 \(\mathrm{B}_q\) 的一个自同构；令 \(\mathrm{A}'\) 为由此得到的 \(\mathrm{A} = \mathrm{Aut}(\mathrm{B}_q)\) 的子群。证明 \(\mathrm{A}' \neq \mathrm{A}\)。（令 \(\varphi\) 为端同态 \(x \mapsto [X_{+}^{\bullet 2}, x]\)；它是 \(\mathrm{B}_q\) 的局部幂零自同态，且 \(e^{\varphi} \in \mathrm{A}, e^{\varphi} \notin \mathrm{A}'\)。）
\item
  群 A 通过结构迁移作用在 \(B_{q}\) 的单表示类集合上。令 \(\pi_{1}\)（分别为 \(\pi_{2},\pi_{3}\)）为习题 13 b) 的 \(\rho_{v,q}\) 型表示（分别为习题 13 的 \(\rho_{v,q}^{\pm}\) 型表示，分别为习题 12 的 \(\rho_{q}\) 型表示）。则 \(A\pi_{1},A\pi_{2}\) 和 \(A\pi_{3}\) 两两不相交。（利用 a)、b)、c)、d)。若 \(\psi\in A\) 使得 \(\psi\pi_{1}\) 是习题 13 b) 的 \(\rho_{v,q}\) 型，则 \(\psi\in A'\)（利用 a) 和 b)）。推出，若 \(\psi\) 是 e) 中的自同构 \(e^{\varphi}\)，则表示 \(\sigma=\pi_{1}\circ\psi\) 具有如下性质：对所有 \(x\in s-\{0\}\)，\(\sigma(x)\) 都没有特征值。 \(^{8}\)
\end{enumerate}

\footnotetext[8]{关于习题 12、13 和 14 的更多细节，参见：D. ARNAL 和 G. PINCZON，《关于 sl(2) 李代数的代数不可约表示》，J. Math. Phys., Vol. 15 (1974), pp. 350-359。}

\begin{enumerate}
\def\labelenumi{\arabic{enumi})}
\setcounter{enumi}{14}
\tightlist
\item
  设 \(\mathfrak{g}\) 是一个 3 维李代数。证明以下条件等价（参见~第 I 章，§6，习题 23）。
\end{enumerate}

\begin{enumerate}
\def\labelenumi{(\roman{enumi})}
\item
  \(\mathfrak{g} = [\mathfrak{g},\mathfrak{g}]\)。
\item
  \(\mathfrak{g}\) 的 Killing 型非退化。
\item
  g 是半单的。
\item
  g 是单的。
\end{enumerate}

¶ 16) 设 g 是一个 3 维单李代数，且 \(\Phi\) 是它的 Killing 型。记 \(\mathfrak{o}(\Phi)\) 为 \(\Phi\) 的正交代数，即由 \(\mathfrak{gl}(\mathfrak{g})\) 中保持 \(\Phi\) 不变的元素组成的子代数。

\begin{enumerate}
\def\labelenumi{\alph{enumi})}
\item
  证明 ad：\(\mathfrak{g} \to \mathfrak{o}(\varPhi)\) 是一个同构。
\item
  证明以下性质等价：
\end{enumerate}

\begin{enumerate}
\def\labelenumi{(\roman{enumi})}
\item
  \(\mathfrak{g}\) 含有一个非零迷向向量（相对于 \(\varPhi\)）。
\item
  \(\mathfrak{g}\) 含有一个非零幂零元。
\item
  g 同构于 s。
\end{enumerate}

\begin{enumerate}
\def\labelenumi{\alph{enumi})}
\setcounter{enumi}{2}
\item
  证明存在一个扩域 \(k_{1}\)，其基域为 \(k\)，次数为 \(\leq 2\)，使得 \(\mathfrak{g}_{(k_1)} = k_1 \otimes_k \mathfrak{g}\) 同构于 \(\mathfrak{s}_{(k_1)}\)。
\item
  在向量空间 \(\mathbf{A} = k\oplus \mathfrak{g}\) 上赋予唯一的代数结构，使 1 为单位元，并使得乘法 \(\mathfrak{g}\times \mathfrak{g}\to \mathrm{A}\) 由公式给出
\end{enumerate}

\[
x. y = \frac {1}{8} \Phi (x, y). 1 + \frac {1}{2} [ x, y ] \quad (x, y \in \mathfrak {g}).
\]

证明 A 是 k 上的四元数代数，而 g 是 A 中由迹为零的元素组成的李子代数（通过扩张基域归约到 g = s 的情形，并证明此时 A 可与矩阵代数 \(\mathbf{M}_{2}(k)\) 认同）。

反之，若 D 是 k 上的四元数代数，则 D 中迹为零的元素构成一个 3 维单李代数，并且相应的代数 A 可与 D 认同。

\begin{enumerate}
\def\labelenumi{\alph{enumi})}
\setcounter{enumi}{4}
\tightlist
\item
  证明下列公式
\end{enumerate}

\[
\begin{array}{c} \Phi (x, x) = - 8 \mathrm{Nrd} _ {\mathrm{A}} (x), \quad \Phi (x, y) = 4 \mathrm{Trd} _ {\mathrm{A}} (x y) \\ 2 [ x, [ y, z ] ] = \Phi (x, y) z - \Phi (x, z) y \qquad (x, y, z \in \mathfrak {g}). \end{array}
\]

\(f\) ) 证明 \(\varPhi\) 的判别式（相对于 \(\mathfrak{g}\) 的任意基）具有 \(-2\lambda^2\) 的形式，其中 \(\lambda \in k^{*}\)。

\begin{enumerate}
\def\labelenumi{\alph{enumi})}
\setcounter{enumi}{6}
\tightlist
\item
  设 n 是一个 \(\geq 0\) 的整数。证明 g-模 \(\mathbf{S}^{n}(\mathfrak{g})\) 有一个唯一的维数为 \(2n + 1\) 的单子模，且此模是绝对单的（归约到 g = s 的情形，并利用习题 9）。
\end{enumerate}

证明 g 只有在同构于 s（即~只有在 A 同构于 \(\mathbf{M}_{2}(k)\)）时才有维数为 2n 的绝对单模。（设 V 是这样的模。如果 \(n \geq 2\)，利用习题 6 c) 证明 g-模 \(V \otimes g\) 有唯一的维数为 2n - 2 的绝对单子模。然后归约到 n = 1 的情形，这显然。）

¶ 17) 保留上一个习题的记号。证明，对所有 \(n \geq 1\)，代数 \(\mathrm{U}(\mathfrak{g})\) 有唯一的双边理想 \(m_{n}\)，使得 \(\mathrm{U}(\mathfrak{g}) / \mathfrak{m}_{n}\) 是维数为 \(n^{2}\) 的中心单代数（扩张标量以归约到 g = s 的情形，并证明此时 \(m_{n}\) 是同态 \(\mathrm{U}(\mathfrak{g}) \to \mathrm{End}(\mathrm{V}(n-1))\) 的核）。\(\mathrm{U}(\mathfrak{g})\) 的每个有限余维双边理想都形如 \(m_{n_{1}} \cap m_{n_{2}} \cap \cdots \cap m_{n_{h}}\)，其中 \(n_{1}, \ldots, n_{h}\) 互异；其余维为 \(n_{1}^{2} + \cdots + n_{h}^{2}\)（应用稠密性定理）。\(m_{n}\) 是 \(\mathrm{U}(\mathfrak{g})\) 中唯一的有限余维极大双边理想。

证明 \(m_{2}\)（作为双边理想）由元素 \(x^{2}-\frac{1}{8}\Phi(x,x)\)（\(x\in g\)）生成，并且商 \(\mathrm{U}(\mathfrak{g})/\mathfrak{m}_{2}\) 可与习题 16 的四元数代数 A 认同。

当 g 同构于 s 时，\(\mathrm{U}(\mathfrak{g})/\mathfrak{m}_{n}\) 同构于 \(\mathbf{M}_{n}(k)\)。利用习题 16 g)，证明当 g 不同构于 s 时，\(\mathrm{U}(\mathfrak{g})/\mathfrak{m}_{n}\) 同构于 \(\mathbf{M}_{n}(k)\) 当且仅当 n 为奇数。\(^{9}\)

\footnotetext[9]{可以证明，当 n 为偶数时，\(\mathrm{U}(\mathfrak{g})/\mathfrak{m}_{n}\) 同构于 \(\mathbf{M}_{n/2}(\mathrm{A})\)。}

¶ 18) 假设 \(k = \mathbf{R}, \mathbf{C}\)，或者 k 是一个关于离散赋值完备且剩余域特征为 \(p \neq 0\) 的域（例如 \(p\)-进域 \(\mathbf{Q}_p\) 的有限扩域）。

\begin{enumerate}
\def\labelenumi{\alph{enumi})}
\tightlist
\item
  设 \(\mathfrak{n}\) 是一个幂零李代数，N 是通过赋予 \(\mathfrak{n}\) Hausdorff 作用律得到的李群（第 III 章，§9，第 5 节），且 \(\rho : \mathfrak{n} \to \mathfrak{gl}(\mathrm{E})\) 是 \(\mathfrak{n}\) 在有限维向量空间 E 上的线性表示。假设 \(\rho(x)\) 对所有 \(x \in \mathfrak{n}\) 都是幂零的，并置 \(\pi(x) = \exp(\rho(x))\)。证明 \(\pi\) 是唯一的李群同态 \(\varphi : \mathrm{N} \to \mathbf{GL}(\mathrm{E})\)，使得 \(\mathrm{L}(\varphi) = \rho\)。
\end{enumerate}

（当 k = R 或 C 时，利用 N 连通这一事实。当 k 是超度量的时，证明 \(\varphi(x)\)（\(x \in N\)）的特征值等于 1；为此，先证明：如果 \(k'\) 是 k 的有限扩域，且 \((\lambda_{1}, \ldots, \lambda_{n}, \ldots)\) 是 \(k'\) 中元素的序列，满足对所有 n 有 \(\lambda_{n} = \lambda_{n+1}^{p}\)，且 \(\lambda_{1} = 1\)，则对所有 n 有 \(\lambda_{n} = 1\)。）

\begin{enumerate}
\def\labelenumi{\alph{enumi})}
\setcounter{enumi}{1}
\tightlist
\item
  设 \(\rho : \mathfrak{s} \to \mathfrak{gl}(\mathrm{E})\) 是 \(\mathfrak{s}\) 的有限维线性表示，并设 \(\pi\) 是从 \(\mathbf{SL}(2, k)\) 到 \(\mathbf{GL}(\mathrm{E})\) 的与 \(\rho\) 相容的同态（第 4 节）。证明 \(\pi\) 是唯一的李群同态 \(\varphi : \mathbf{SL}(2, k) \to \mathbf{GL}(\mathrm{E})\)，使得 \(\mathrm{L}(\varphi) = \rho\)。（利用 \(a\) ) 证明 \(\pi\) 与 \(\rho\) 在 \(\exp(n)\) 上一致，其中 \(n\) 是 \(\mathfrak{s}\) 中的幂零元，并注意到 \(\mathbf{SL}(2, k)\) 由 \(\exp(n)\) 生成。）
\end{enumerate}

\exercisesection{\S~2}

\begin{enumerate}
\def\labelenumi{\arabic{enumi})}
\tightlist
\item
  设 \(\mathfrak{g}\) 是一个 3 维单李代数，且 \(\varPhi\) 是其 Killing 型（参见~§1，习题 15、16、17）。
\end{enumerate}

\begin{enumerate}
\def\labelenumi{\alph{enumi})}
\item
  元素 \(x \in g\) 是正则元当且仅当 \(\Phi(x, x) \neq 0\)。令 \(h_x = kx\) 为由此类元素生成的嘉当子代数。证明 \(h_x\) 是分裂的当且仅当 \(2\Phi(x, x)\) 是 k 中的平方。
\item
  证明
\end{enumerate}

\(\mathfrak{g}\) 可分裂 \(\Longleftrightarrow \mathfrak{g}\) 同构于 \(\mathfrak{sl}(2,k)\) 。

\begin{enumerate}
\def\labelenumi{\arabic{enumi})}
\setcounter{enumi}{1}
\tightlist
\item
  设 \(k_{1}\) 是 \(k\) 的一个有限扩域，次数 \(n \geq 2\)。
\end{enumerate}

\begin{enumerate}
\def\labelenumi{\alph{enumi})}
\item
  证明半单 k-代数 \(\mathfrak{sl}(2,k_{1})\) 不是可分裂的。
\item
  证明可分裂的单 k-代数 \(\mathfrak{sl}(n,k)\) 包含一个非分裂的嘉当子代数 \(h_{1}\)。（选择代数 \(k_{1}\) 到 \(\mathbf{M}_{n}(k)\) 的嵌入，并取 \(\mathfrak{h}_{1}=k_{1}\cap\mathfrak{sl}(n,k)\)。）
\end{enumerate}

\begin{enumerate}
\def\labelenumi{\arabic{enumi})}
\setcounter{enumi}{2}
\tightlist
\item
  设 \((\mathfrak{g},\mathfrak{h})\) 是一个分裂半单李代数，R 是其根系，K 是 \(\mathfrak{g}\) 的 Killing 型在 \(\mathfrak{h}\) 上的限制。采用第 VI 章 §1 第 12 节的记号，有 \(K = B_{R^{\vee}}\)，且当 R 不可约时有 \(K = 4\gamma(R)\varPhi_{R^{\vee}}\)；此外，如果 R 的所有根长度相同，则
\end{enumerate}

\[
\mathrm{K} (H _ {\alpha}, H _ {\alpha}) = 4 h \quad \text { for   all } \alpha \in \mathrm{R},
\]

其中 \(h\) 是 \(\mathrm{W}(\mathbb{R})\) 的 Coxeter 数。

例如，如果 \(\mathfrak{g}\) 是 \(\mathrm{E}_6, \mathrm{E}_7\) 或 \(\mathrm{E}_8\) 型的单李代数，则 \(\mathrm{K}(H_{\alpha}, H_{\alpha})\) 等于 48、72 或 120。

\begin{enumerate}
\def\labelenumi{\arabic{enumi})}
\setcounter{enumi}{3}
\tightlist
\item
  设 \((\mathfrak{g},\mathfrak{h})\) 是一个分裂半单李代数，且 \((X_{\alpha})_{\alpha \in \mathrm{R}}\) 是满足引理 2 条件的元素族。如果 \(\alpha ,\beta \in \mathrm{R}\) 且 \(\alpha +\beta \in \mathrm{R}\)，则定义 \(\mathrm{N}_{\alpha \beta}\)，使其满足公式 \([X_{\alpha},X_{\beta}] = \mathrm{N}_{\alpha \beta}X_{\alpha +\beta}\)；如果 \(\alpha +\beta \notin \mathrm{R}\)，则置 \(\mathrm{N}_{\alpha \beta} = 0\)。证明下列公式：
\end{enumerate}

\(a)\mathrm{N}_{\alpha \beta} = -\mathrm{N}_{\beta \alpha}.\)

\begin{enumerate}
\def\labelenumi{\alph{enumi})}
\setcounter{enumi}{1}
\tightlist
\item
  如果 \(\alpha, \beta, \gamma \in \mathbb{R}\) 满足 \(\alpha + \beta + \gamma = 0\)，则
\end{enumerate}

\[
\frac {\mathrm{N} _ {\alpha \beta}}{\langle \gamma , \gamma \rangle} = \frac {\mathrm{N} _ {\beta \gamma}}{\langle \alpha , \alpha \rangle} = \frac {\mathrm{N} _ {\gamma \alpha}}{\langle \beta , \beta \rangle}.
\]

\begin{enumerate}
\def\labelenumi{\alph{enumi})}
\setcounter{enumi}{2}
\tightlist
\item
  如果 \(\alpha, \beta, \gamma, \delta \in \mathbb{R}\) 满足 \(\alpha + \beta + \gamma + \delta = 0\)，且没有一对元素的和为零，则有：
\end{enumerate}

\[
\frac {\mathrm{N} _ {\alpha \beta} \mathrm{N} _ {\gamma \delta}}{\langle \gamma + \delta , \gamma + \delta \rangle} + \frac {\mathrm{N} _ {\beta \gamma} \mathrm{N} _ {\alpha \delta}}{\langle \alpha + \delta , \alpha + \delta \rangle} + \frac {\mathrm{N} _ {\gamma \alpha} \mathrm{N} _ {\beta \delta}}{\langle \beta + \delta , \beta + \delta \rangle} = 0.
\]

\begin{enumerate}
\def\labelenumi{\arabic{enumi})}
\setcounter{enumi}{4}
\item
  设 \((\mathfrak{g},\mathfrak{h})\) 是一个分裂半单李代数，且 \((X_{\alpha})_{\alpha\in\mathbb{R}}\) 是满足引理 2 条件的元素族。令 B 为 R 的一个基。\(H_{\alpha}\)（\(\alpha\in B\)）和 \(X_{\alpha}\)（\(\alpha\in R\)）构成向量空间 g 的一个基。证明相对于此基（《代数》第 IX 章，§2），g 的 Killing 型的判别式是一个有理数，与 h、B 和 \(X_{\alpha}\) 的选择无关。¹⁰ 推出，如果 \(n=\dim g\)，则 \(\bigwedge^{n}g\) 中由 \(H_{\alpha}\)（\(\alpha\in B\)）和 \(X_{\alpha}\)（\(\alpha\in R\)）的外积定义的元素，除差一个符号外，与 h、B 和 \(X_{\alpha}\) 的选择无关。

\footnotetext[10]{关于这个判别式的显式计算，参见：T. A. SPRINGER 和 R. STEINBERG，《共轭类》（第 4.8 节），《代数群及相关有限群研讨会》，数学讲义 131，Springer-Verlag（1970）。}
  \item
  设 \((X_{\alpha})_{\alpha\in\mathbb{R}}\) 是分裂半单李代数 \((\mathfrak{g},\mathfrak{h})\) 中的一个谢瓦莱系。令 \(\alpha,\beta\in R\)，并令 p（分别为 q）是满足 \(\beta+j\alpha\in R\)（分别为 \(\beta-j\alpha\in R\)）的最大整数 j，参见~引理 4。证明
\end{enumerate}

\[
\operatorname{ad} (X _ {\alpha}) ^ {k} (X _ {\beta - q \alpha}) = \pm k! X _ {\beta + (k - q) \alpha} \text { for } 0 \leq k \leq p + q.
\]

推出命题 8 中的代数 \(\mathfrak{g}_{\mathbf{Z}}\) 在 \(\mathrm{ad}(X_{\alpha})^k / k!\) 和 \(e^{\mathrm{ad}(X_{\alpha})}\) 下稳定（参见~§12）。

\begin{enumerate}
\def\labelenumi{\arabic{enumi})}
\setcounter{enumi}{6}
\tightlist
\item
  假设 \(k\) 是有序域。设 \((\mathfrak{g},\mathfrak{h})\) 是秩为 \(l\) 的分裂半单李代数；置 \(\dim \mathfrak{g} = l + 2m\)。
\end{enumerate}

\begin{enumerate}
\def\labelenumi{\alph{enumi})}
\item
  证明 \(\varPhi\)（g 的 Killing 型）是一个秩为 2m 的中性型与一个秩为 l 的正非退化型的直和；特别地，其指标为 m。
\item
  设 \(\varphi\) 是 \(\mathfrak{g}\) 的一个对合自同构，其在 \(\mathfrak{h}\) 上的限制为 -Id。证明下列型
\end{enumerate}

\[
(x, y) \mapsto - \Phi (x, \varphi (y)), \quad x, y \in \mathfrak {g},
\]

是对称的、非退化的且正的。

\(c)\) 令 \(k' = k(\sqrt{\alpha})\)，其中 \(\alpha\) 是一个 \(< 0\) 的元素，属于 \(k\)。记 \(c\) 为 \(k\) 上 \(k'\) 的非平凡自同构。令 \(\mathfrak{g}_c\) 为 \(k\)-子空间 \(\mathfrak{g}_{(k')} = k' \otimes_k \mathfrak{g}\) 中由满足下式的元素 \(y\) 构成：

\[
(1 \otimes \varphi). y = (c \otimes 1). y.
\]

证明 \(g_{c}\) 是 \(\mathfrak{g}_{(k^{\prime})}\) 的 k-李子代数，且 \(g_{c}\) 到 \(\mathfrak{g}_{(k^{\prime})}\) 的嵌入延拓为从 \(k^{\prime}\otimes_{k}g_{c}\) 到 \(\mathfrak{g}_{(k^{\prime})}\) 的同构。代数 \(g_{c}\) 是半单的，而 \(\sqrt{\alpha}h\) 是它的一个嘉当子代数。

\(d\) ) 证明 \(\mathfrak{g}_c\) 的 Killing 型是负定的。推出 \(\mathfrak{g}_c\) 不是可分裂的（除非 \(\mathfrak{g} = 0\)）。

\begin{enumerate}
\def\labelenumi{\alph{enumi})}
\setcounter{enumi}{4}
\tightlist
\item
  当 \(k = \mathbf{R}\) 时，证明 \(\operatorname{Int}(\mathfrak{g}_c)\) 是紧的。
\end{enumerate}

\begin{enumerate}
\def\labelenumi{\arabic{enumi})}
\setcounter{enumi}{7}
\tightlist
\item
  \begin{enumerate}
  \def\labelenumii{\alph{enumii})}
  \tightlist
  \item
    设 g 是李代数，n 为一个 \(\geq 0\) 的整数。令 \(\Sigma_{n}(\mathfrak{g})\) 为 \(g^{n}\) 的子集，由作为 k-代数生成 g 的族 \((x_{1},\ldots,x_{n})\) 组成。证明 \(\Sigma_{n}(\mathfrak{g})\) 在 Zariski 拓扑中是 \(g^{n}\) 的开集（第 VII 章附录 I）。如果 \(k'\) 是 k 的扩域，则 \(\Sigma_{n}(\mathfrak{g}_{(k')})\cap\mathfrak{g}^{n}=\Sigma_{n}(\mathfrak{g})\)。推出，如果 \(\mathfrak{g}_{(k')}\) 可由 n 个元素生成，则 g 也可以。
  \end{enumerate}
\end{enumerate}

\begin{enumerate}
\def\labelenumi{\alph{enumi})}
\setcounter{enumi}{1}
\item
  设 \((\mathfrak{g},\mathfrak{h})\) 是一个分裂半单李代数。令 \(x\) 是 \(\mathfrak{h}\) 中满足：\(\alpha (x)\neq 0\) 对所有 \(\alpha \in \mathrm{R}\) 成立，且 \(\alpha (x)\neq \beta (x)\) 对任意两个不同的 \(\alpha ,\beta \in \mathrm{R}\) 成立的元素。对所有 \(\alpha \in \mathrm{R}\)，令 \(y_{\alpha}\) 为 \(\mathfrak{g}^{\alpha}\) 的一个非零元素，并令 \(y = \sum_{\alpha \in \mathrm{R}}y_{\alpha}\)。证明，对所有 \(\alpha \in \mathrm{R}\)，存在一个无常数项的多项式 \(\mathrm{P}_{\alpha}(\mathrm{T})\in k[\mathrm{T}]\)，使得 \(y_{\alpha} = \mathrm{P}_{\alpha}(\mathrm{ad}x).y\)。推出 \(\mathfrak{g}\) 由 \(\{x,y\}\) 生成。
\item
  利用 \(a\) ) 和 \(b\) )，证明每个半单李代数都可由两个元素生成。
\end{enumerate}

¶9) 设 G 是一个连通有限维实李群。设 g 为其李代数，且 \(\{x_{1},\ldots,x_{n}\}\) 是 g 的一个生成族。对 \(m\geq0\)，记 \(\Gamma_{m}\) 为由 \(\exp(2^{-m}x_{i})\)（\(1\leq i\leq n\)）生成的 G 的子群。我们有 \(\Gamma_{0}\subset\Gamma_{1}\subset\cdots\)。

\begin{enumerate}
\def\labelenumi{\alph{enumi})}
\item
  证明 \(\Gamma_{m}\) 的并在 G 中稠密。
\item
  令 \(H_{m}\) 为 \(\overline{\Gamma}_{m}\)（\(\Gamma_{m}\) 的闭包）的单位元连通分支。我们有 \(H_{0} \subset H_{1} \subset \cdots\)，且族（\(H_{m}\)）是稳定的；令 H 为 m 充分大时 \(H_{m}\) 的公共值。证明 H 在 G 中正规（注意 H 被所有 \(\Gamma_{m}\) 规范化），且 \(\Gamma_{m}\) 在 G/H 中的像是 G/H 的离散子群。这些子群的并在 G/H 中稠密，推出（参见~第 III 章，§6，习题 23 d)）G/H 是幂零的。
\item
  假设 \(\mathfrak{g} = \mathcal{D}\mathfrak{g}\)。证明 \(\mathrm{G} = \mathrm{H}\)，换言之，\(\Gamma_{m}\) 在 \(\mathrm{G}\) 中稠密，当 \(m\) 充分大时。
\end{enumerate}

\begin{enumerate}
\def\labelenumi{\arabic{enumi})}
\setcounter{enumi}{9}
\item
  设 \(G\) 是一个连通半单实李群。利用习题 8 和 9 证明，存在一个由两个元素生成的 \(G\) 的稠密子群。
\item
  设 \(\mathrm{R}\) 是秩为 \(l\) 的根系，\(\varPhi_{\mathrm{R}}\) 是相应的典则双线性型（第 VI 章，§1，第 12 节）。矩阵 \(\varPhi = (\varPhi_{\mathrm{R}}(\alpha, \beta))_{\alpha, \beta \in \mathrm{R}}\) 的秩为 \(l\)，且 \(\varPhi^2 = \varPhi\)。推出公式 \(\sum_{\alpha \in \mathrm{R}} \varPhi_{\mathrm{R}}(\alpha, \alpha) = \operatorname{Tr} \varPhi = l\)。
\end{enumerate}

\exercisesection{\S~3}

\begin{enumerate}
\def\labelenumi{\arabic{enumi})}
\item
  设 \(\Gamma\) 是 \(\mathrm{Q}(\mathrm{R})\) 的有限指数子群，且 \(\mathrm{P} = \Gamma \cap \mathrm{R}\)。代数 \(\mathfrak{h} + \mathfrak{g}^{\mathrm{P}}\) 是约化的，并且 \(\mathfrak{g}\) 中任何包含 \(\mathfrak{h}\) 的约化子代数都可以用这种方式得到（利用第 VI 章，§1，习题 6 b)）。
\item
  令 X 为 g 的约化子代数集合中与 g 不同且包含 h 的那些。利用第 VI 章，§4，习题 4 确定 X 的极小元素。证明这种极大子代数的中心维数为 0 或 1，分别对应于所讨论习题的 a) 或 b) 情形。
\item
  设 \(\mathfrak{b}\) 是 \(\mathfrak{g}\) 的一个 Borel 子代数，且 \(l = \mathrm{rk}(\mathfrak{g})\)。证明代数 \(\mathfrak{b}\) 的最少生成元数是 \(l\)，当 \(l \neq 1\) 时，且在 \(l = 1\) 时为 2。
\item
  假设 \(k = \mathbf{R}\) 或 \(\mathbf{C}\)。置 \(G = \operatorname{Int}(\mathfrak{g})\)，并将 \(G\) 的李代数与 \(\mathfrak{g}\) 认同。设 \(\mathfrak{b}\) 是 \((\mathfrak{g},\mathfrak{h})\) 的一个 Borel 子代数，且 \(\mathfrak{n}\) 是其幂零元素的集合。记 \(H,B,N\) 为 \(G\) 中李代数分别为 \(\mathfrak{h},\mathfrak{b},\mathfrak{n}\) 的积分子群。证明 \(H,B,N\) 是 \(G\) 的李子群，\(N\) 是单连通的，且 \(B\) 是 \(H\) 与 \(N\) 的半直积。
\item
  设 \(\mathfrak{m}\) 是半单李代数 \(\mathfrak{a}\) 的一个抛物子代数。
\end{enumerate}

\begin{enumerate}
\def\labelenumi{\alph{enumi})}
\item
  设 p 是 m 的子代数。证明 p 是 a 的抛物子代数当且仅当 p 包含 m 的根基 r，且 p/r 是半单代数 m/r 的抛物子代数。
\item
  若 \(m'\) 是 a 的一个抛物子代数，则 \(m \cap m'\) 的每个嘉当子代数都是 a 的嘉当子代数。（归约到分裂情形，并将命题 10 应用于包含在 m 和 \(m'\) 中的 Borel 子代数。）
\end{enumerate}

\begin{enumerate}
\def\labelenumi{\arabic{enumi})}
\setcounter{enumi}{5}
\item
  半单李代数 \(\mathfrak{a}\) 的两个 Borel 子代数若其交是嘉当子代数，则称它们相反。证明，如果 \(\mathfrak{b}\) 是 \(\mathfrak{a}\) 的一个 Borel 子代数，且 \(\mathfrak{h}\) 是 \(\mathfrak{b}\) 的一个嘉当子代数，则存在唯一的 \(\mathfrak{a}\) 的 Borel 子代数，它与 \(\mathfrak{b}\) 相反且包含 \(\mathfrak{h}\)。（归约到分裂情形。）
\item
  设 \(\mathfrak{a}\) 是半单李代数，\(\mathfrak{h}\) 是 \(\mathfrak{a}\) 的嘉当子代数，且 \(\mathfrak{s}\) 是 \(\mathfrak{a}\) 的一个包含 \(\mathfrak{h}\) 的半单子代数。证明：
\end{enumerate}

\[
(\mathfrak {a}, \mathfrak {h}) \text {is split} \Longleftrightarrow (\mathfrak {s}, \mathfrak {h}) \text {is split.}
\]

构造一个例子，其中 \(\mathfrak{a}\) 可分裂但 \(\mathfrak{s}\) 不可分裂。（取 \(\mathfrak{a} = \mathfrak{sp}(4, k)\) 和 \(\mathfrak{s} = \mathfrak{sl}(2, k')\)，其中 \(k'\) 是 \(k\) 的二次扩域。）

\begin{enumerate}
\def\labelenumi{\arabic{enumi})}
\setcounter{enumi}{7}
\tightlist
\item
  选择 R 的一个基，因而得到一组正根 \(R_{+}\)。假设 R 不可约。令 \(\tilde{\alpha}\) 为最高根，S 为与 \(\tilde{\alpha}\) 正交的根的集合，且 \(S_{+} = S \cap R_{+}\)。
\end{enumerate}

\begin{enumerate}
\def\labelenumi{\alph{enumi})}
\tightlist
\item
  令
\end{enumerate}

\[
\mathfrak {g} ^ {\prime} = \mathfrak {g} ^ {\mathrm{S}} \oplus \mathfrak {h} _ {\mathrm{S}}, \quad \mathfrak {m} _ {+} = \sum_ {\alpha \in \mathrm{S} _ {+}} \mathfrak {g} ^ {\alpha}, \quad \mathfrak {m} _ {-} = \sum_ {\alpha \in \mathrm{S} _ {+}} \mathfrak {g} ^ {- \alpha}.
\]

则 \(\mathfrak{g}'\) 是半单的且 \(\mathfrak{g}' = \mathfrak{h}_{\mathrm{S}} \oplus \mathfrak{m}_{+} \oplus \mathfrak{m}_{-}\)。

\begin{enumerate}
\def\labelenumi{\alph{enumi})}
\setcounter{enumi}{1}
\tightlist
\item
  置：
\end{enumerate}

\[
\mathfrak {n} _ {+} = \sum_ {\alpha \in \mathrm{R} _ {+}} \mathfrak {g} ^ {\alpha}, \quad \mathfrak {p} = \sum_ {\alpha \in \mathrm{R} _ {+} - \mathrm{S} _ {+}} \mathfrak {g} ^ {\alpha}, \quad \mathfrak {p} _ {0} = \sum_ {\alpha \in \mathrm{R} _ {+} - \mathrm{S} _ {+} - \{\tilde {\alpha} \}} \mathfrak {g} ^ {\alpha}.
\]

则

\[
\mathfrak {n} _ {+} = \mathfrak {m} _ {+} \oplus \mathfrak {p}, \quad [ \mathfrak {m} _ {+}, \mathfrak {p} _ {0} ] \subset \mathfrak {p} _ {0}, \quad [ \mathfrak {n} _ {+}, \mathfrak {g} ^ {\tilde {\alpha}} ] = 0.
\]

特别地，\(\mathfrak{p}\) 是 \(\mathfrak{n}_{+}\) 的理想。

\begin{enumerate}
\def\labelenumi{\alph{enumi})}
\setcounter{enumi}{2}
\tightlist
\item
  对所有 \(\alpha \in \mathrm{R}_{+} - \mathrm{S}_{+} - \{\tilde{\alpha}\}\)，存在唯一的 \(\alpha' \in \mathrm{R}_{+} - \mathrm{S}_{+} - \{\tilde{\alpha}\}\)，使得 \(\alpha + \alpha' = \tilde{\alpha}\)。对所有 \(\alpha \in \mathrm{R}_{+} - \mathrm{S}_{+} - \{\tilde{\alpha}\}\)，可以选择非零元素 \(X_{\alpha}\)，它属于 \(\mathfrak{g}^{\alpha}\)，使得
\end{enumerate}

\[
\begin{array}{r l} {[ X _ {\alpha}, X _ {\alpha^ {\prime}} ] = \pm X _ {\tilde {\alpha}}} & {\mathrm{if} \alpha \in \mathrm{R} _ {+} - \mathrm{S} _ {+} - \{\tilde {\alpha} \}} \\ {[ X _ {\alpha}, X _ {\beta} ] = 0} & {\mathrm{if} \alpha , \beta \in \mathrm{R} _ {+} - \mathrm{S} _ {+}, \beta \neq \alpha^ {\prime}. ^ {1 1}} \end{array}
\]

\footnotetext[11]{本习题由 A. JOSEPH 提供给我们。}

\begin{enumerate}
\def\labelenumi{\arabic{enumi})}
\setcounter{enumi}{8}
\tightlist
\item
  构造半单李代数 \(\mathfrak{a}\) 的例子，使得：
\end{enumerate}

\begin{enumerate}
\def\labelenumi{\roman{enumi})}
\item
  a 没有 Borel 子代数；
\item
  a 有一个 Borel 子代数，但不是可分裂的。
\end{enumerate}

¶10) 设 a 是半单李代数，且 x 是 a 的一个元素。如果 ad x 可对角化（《代数》第 VII 章），换言之，如果存在一个由 ad x 的特征向量组成的 a 的基，则称 x 是可对角化的。

\begin{enumerate}
\def\labelenumi{\alph{enumi})}
\item
  设 \(\mathfrak{c}\) 是 \(\mathfrak{a}\) 的一个交换子代数，由可对角化元素组成，且 L 是 \(\mathfrak{c}\) 在表示 \(\mathrm{ad}_{\mathfrak{a}}\) 中的权集。集合 L 是 \(\mathfrak{c}^*\) 的一个有限子集，包含 0（除非 \(\mathfrak{a} = 0\)），且 \(\mathfrak{a} = \bigoplus_{\lambda \in L} \mathfrak{a}^{\lambda}(\mathfrak{c})\)。证明存在 L 的子集 M，使得 \(L - \{0\}\) 是 M 与 \(-M\) 的不交并，且 \((M + M) \cap L \subset M\)。如果 M 具有这些性质，置 \(\mathfrak{a}^{M} = \bigoplus_{\lambda \in M} \mathfrak{a}^{\lambda}(\mathfrak{c})\)，并置 \(\mathfrak{p}^{M} = \mathfrak{a}^{0}(\mathfrak{c}) \oplus \mathfrak{a}^{M}\)。代数 \(\mathfrak{a}^{0}(\mathfrak{c})\) 是 \(\mathfrak{c}\) 在 \(\mathfrak{a}\) 中的交换子；它在 \(\mathfrak{a}\) 中是约化的（第 VII 章，§1，第 5 节），且其嘉当子代数就是 \(\mathfrak{a}\) 的嘉当子代数（第 VII 章，§2，第 3 节）。证明 \(\mathfrak{p}^{M}\) 是 a 的抛物子代数，且 \(a^{M}\) 是 \(p^{M}\) 的根基的幂零元素集合。（利用 c 包含于 a 的一个嘉当子代数这一事实，并通过扩张标量归约到这个子代数是分裂的情形。）
\item
  保留 a) 的假设和记号，并进一步假设 a 可分裂。证明 c 包含于 a 的一个分裂嘉当子代数中。（取 a 的一个包含于 \(p^{M}\) 中的分裂嘉当子代数 h；则存在唯一的 \(h'\)，它是 \(\mathfrak{a}^{0}(\mathfrak{c})\) 的嘉当子代数，使 h 包含于 \(h' \oplus a^{M}\) 中；证明 \(h'\) 是 a 的分裂嘉当子代数，且 \(h'\) 包含 c。）
\end{enumerate}

\begin{enumerate}
\def\labelenumi{\arabic{enumi})}
\setcounter{enumi}{10}
\tightlist
\item
  设 \(\mathfrak{a} = \prod \mathfrak{a}_i\) 是半单李代数的有限积。\(\mathfrak{q}\) 是 \(\mathfrak{a}\) 的子代数，且它是抛物（分别为 Borel）子代数，当且仅当它形如 \(\mathfrak{q} = \prod \mathfrak{q}_i\)，其中对于所有 \(i\)，\(\mathfrak{q}_i\) 是 \(\mathfrak{a}_i\) 的抛物（分别为 Borel）子代数。
\end{enumerate}

¶ 12) 设 \(k'\) 是 \(k\) 的有限扩域，\(\mathfrak{a}'\) 是半单 \(k'\)-李代数，且 \(\mathfrak{a}\) 是其底层 \(k\)-李代数。证明 \(\mathfrak{a}\) 的抛物（分别为 Borel）子代数与 \(\mathfrak{a}'\) 的相同。（将标量扩张到 \(k\) 的代数闭包，并利用习题 11。）

\begin{enumerate}
\def\labelenumi{\arabic{enumi})}
\setcounter{enumi}{12}
\item
  设 p 和 q 是半单李代数 a 的两个抛物子代数，且 n 是 p 的幂零根基。证明 \(\mathfrak{m} = (\mathfrak{p} \cap \mathfrak{q}) + \mathfrak{n}\) 是 a 的抛物子代数。（归约到 a 分裂且 q 为 Borel 子代数的情形；选择包含于 \(p \cap q\) 中的嘉当子代数 h，并确定相应根系的子集 P，使 \(m = h + g^{P}\)。）
\item
保留命题 9 的记号。
\end{enumerate}

\begin{enumerate}
\def\labelenumi{\alph{enumi})}
\item
  令 \(\alpha \in \mathrm{B}\)。证明 \(\mathfrak{n} \cap \operatorname{Kerad} X_{\alpha}\) 是各 \(\mathfrak{g}^{\beta}\) 的直和，其中 \(\beta\) 属于 \(\mathrm{R}_{+}\) 中满足 \(\alpha + \beta \notin \mathrm{R}_{+}\) 的元素集合。
\item
  推出，如果 g 是单的，则 n 的中心等于 \(g^{\tilde{\alpha}}\)，其中 \(\tilde{\alpha}\) 是 \(R_{+}\) 的最高根。在一般情形下，n 的中心维数等于 g 的单分量数。
\end{enumerate}